% ============================================================================
% sec6_experiments.tex — Experiments (Plan B, 2026-11-06)
% Unified-recipe control experiment + run-level reproducibility + LightRR
% motivation + frozen held-out protocol + negative results.
% All numbers are from the Lead's final unified-script table (15-scene panel,
% hand-annotated Omega, uint8 evaluation; three seeds 101/202/303).
% ============================================================================
\section{Experiments}
\label{sec:exp}

\subsection{Protocol and baselines}
\label{sec:setup}
All scores below are on the 15-scene development panel, produced by a frozen
uint8 evaluation script. Trained arms use seeds $101$, $202$ and $303$; we
report the mean and standard deviation over scenes, and differences are
scene-level paired means with $95\%$ confidence intervals over the 15 scenes.
The fine-tuned baseline is a strongly fine-tuned RDNet~\cite{zhao2025rdnet}
(frozen FocalNet backbone, $65.95$M trainable parameters, $6000$ iterations,
$30/40/30$ real/synthetic/public mixture; supplementary material S3). Zero-shot
public methods range from $18.12$\,dB (IBCLN~\cite{li2020ibcln}) to
$20.50$\,dB (Zhu et al.~\cite{zhu2024rrw}), while a trivial identity
predictor scores $18.44$\,dB and the strongly fine-tuned baseline scores
$22.38\pm0.23$\,dB. We compare public methods zero-shot only and make no
public-benchmark state-of-the-art claim.

\subsection{A unified-recipe control experiment}
\label{sec:control}
To ask what a training-time regularizer can achieve on this benchmark, the
regularizer and four controls share one recipe---identical data mixture,
schedule, seeds and evaluation---so only the extra term differs. The arms are:
\emph{base} (no extra term); \emph{identity} (a plain identity penalty
$\lambda\|\That-I\|_1$ with $\lambda=0.409$); \emph{identity-m} (the same with
$\lambda=0.085$, magnitude-matched to the change-aware weight); \emph{hard}
(the change-aware statistic thresholded at $\tau=0.05$); \emph{shuffled} (the
change-aware weights permuted across pixels); and \emph{full} ($\quiet$ with
$\lambda=1$, $\tau=0.05$). Table~\ref{tab:control} reports the result.

\begin{table}[t]
\centering
\footnotesize
\caption{Unified-recipe control experiment (development panel, $15$ scenes,
three seeds). PSNR is the mean over scenes; $\Delta$ is the scene-level paired
mean against \emph{full} with its $95\%$ CI. All regularized arms sit within
$0.44$\,dB of base.}
\label{tab:control}
\begin{tabular}{lcc}
\toprule
arm & PSNR (dB) & $\Delta$ vs.\ full (dB) \\
\midrule
base (no regularizer)        & $22.38\pm0.23$ & $+0.19\,[+0.05,+0.33]$ \\
identity ($\lambda{=}.409$)  & $22.49\pm0.34$ & $+0.079\,[+0.021,+0.137]$ \\
identity-m ($\lambda{=}.085$) & $22.59\pm0.21$ & $-0.014\,[-0.215,+0.186]$ \\
hard ($\tau{=}.05$)          & $22.82\pm0.36$ & $-0.246\,[-0.355,-0.138]$ \\
shuffled weight              & $22.36\pm0.39$ & $+0.209\,[+0.033,+0.386]$ \\
\textbf{full ($\lambda{=}1$, $\tau{=}.05$)} & $\mathbf{22.57\pm0.00}$ & --- \\
\bottomrule
\end{tabular}
\end{table}

The regularized arms span $22.36$--$22.82$\,dB, from $0.02$\,dB below to
$0.44$\,dB above base, and the ordering does not favour the proposed shape: a
magnitude-matched identity penalty reaches $22.59$\,dB, a hard threshold of the
same statistic reaches $22.82$\,dB, and full reaches $22.57$\,dB. The paired
difference full$-$identity is $+0.079$\,dB ($95\%$ CI $[0.021,0.137]$,
$p=0.011$, positive on $14/15$ scenes), but it sits at the contrast's minimum
detectable effect ($0.075$\,dB) and far below the run-level spread of
Sec.~\ref{sec:runlevel}. Full is not better than identity-m
($-0.014$\,dB, CI $[-0.215,0.186]$, $p=0.88$) and is significantly worse than
hard ($-0.246$\,dB, CI $[-0.355,-0.138]$, positive on $2/15$ scenes). Relative
to base, the full arm gains $0.19$\,dB (CI $[0.05,0.33]$, $11/15$ scenes).
The control therefore bounds what this family of regularizers achieves at a
few tenths of a decibel, while showing that the specific change-aware shape is
not the lever.

\subsection{What the control must exceed: run-level reproducibility}
\label{sec:runlevel}
Seed spread on this panel is below $0.44$\,dB, so a naive reading would treat
$0.08$\,dB as measurable. It is not. Re-running the same arm with the same
script and seed changed the panel score by $-0.88$\,dB for base
($22.998\rightarrow22.119$) and by $-0.51$\,dB for full
($23.079\rightarrow22.570$); cross-date differences reached $+1.008$\,dB for
base and $+0.369$\,dB for full. Run-level uncertainty is therefore at least as
large as seed-level uncertainty, and no difference of the size discussed above
can be read as a property of a regularizer. We report these re-runs as part of
the protocol and use them to set the resolution at which any claim on this
panel is made. Before the recipes were unified, an earlier comparison of the
two training scripts produced a $+0.50$\,dB difference; because the scripts
differ in recipe, that number is a historical observation only and is not a
result.

\subsection{In-domain fine-tuning does not rescue a lightweight method}
\label{sec:lightrr}
One might assume that the benchmark is hard because the model is
under-adapted. It is not. We fine-tune LightRR~\cite{yin2026lightrr} in-domain
with its official recipe for $102{,}500$ iterations ($41$ epochs; one more than
pre-registered, disclosed). The best of two learning-rate runs reaches
$20.221$\,dB ($20.221$ at $10^{-4}$ and $20.198$ at $5\times10^{-5}$), more
than $2$\,dB below the fine-tuned baseline, and it wins on no scene ($0/15$).
The residual-reflection coefficient tells the same story: it moves from
$-0.859$ zero-shot to only $-0.489$/$-0.511$ after in-domain fine-tuning,
whereas the strongly fine-tuned RDNet reaches $-0.29$ to $-0.34$. The deficit
lies in what this data protocol asks of the method, not in the amount of
in-domain exposure.

\subsection{A frozen protocol for an independent capture}
\label{sec:heldout}
The panel was used during development, so it cannot confirm. We therefore
pre-registered a one-shot protocol on a new capture through the same film
(indoor and outdoor scenes, disjoint from the development sessions): the model,
the baselines, the evaluation script and the primary contrast are frozen before
any held-out frame is scored, no tuning is performed on the held-out set, and
per-scene results are reported with the same paired intervals. If the primary
contrast crosses zero, the claim is reduced rather than re-tuned. At submission
time the capture is still being collected; we report the protocol, make no
held-out number, and will release the frozen scores with the data.

\subsection{Negative results}
\label{sec:negative}
Several extensions failed and we report them rather than drop them:
perceptual-loss weighting hedged against the structural gap; splicing the
input's high frequencies back into the output restored detail but lowered PSNR;
mixing occluder-region data raised region-restricted PSNR while degrading the
residual-reflection coefficient; photometric adapters and gates were rejected
by cheap probes; and averaging three independently trained models failed
because the seeds converge to different basins. Each is documented in the
supplementary material.
